% ============================================================
% 纯答案册·课时08-直线与圆的位置关系 —— 工具/答案抽册器.py 抽册生成件（勿手改；第三档＝纯答案单独成册）
% 源件（只读）：C:/提示词/工作区/M3-第2章量产0913/成卷/导学件/课时08-直线与圆的位置关系/main.tex（md5 8e7ac74c854066b08710cee7ebe28778）
%% 口径：题面不重印；逐题＝题号(印面号同正册)＋[答案]＋[详解]，值/详解照源件逐字
%       （钉值门硬断言：册值≡正册件内值≡值台账值tex，逐字全等）。册序＝正册装配序＝印面号 1..21 升序（同序同号）；键序断言基准＝题面库冻结 manifest canonical 键序。
%% 三档导出：整体 main-true／纯题 main-false（在产）｜本册＝纯答案册（本件）
%% 双档：ansbook-true.tex＝详解印本；ansbook-false.tex＝纯值速查本（详解不印）
% 编译：xelatex <壳> ×2，cwd＝本目录；sty＝qp-m3.sty 本地挂载（md5 7c3930362be8a0a2bdf21bbf8ac16573，锁校验）
% ============================================================
\documentclass[fontset=none]{ctexart}
\usepackage{qp-m3}
% —— 印面逐字符号路由补钉（承源件；− 走 TNR、▱ NSC 子块；禁改 sty 保 md5） ——
\xeCJKDeclareCharClass{Default}{"2212}%
\setCJKfamilyfont{nscblk}[Path=C:/提示词/工作区/字替对照-0909/variantF/fonts/,BoldFont=NSC-w600.ttf]{NSC-w500.ttf}%
\xeCJKDeclareSubCJKBlock{nscblk}{"25B1}%
\newcommand{\pxparallelogram}{{\CJKfamily{nscblk}▱}}%
% —— 断行微调（承母版实测档，三零门承重；\emergencystretch 不另设） ——
\xeCJKsetup{CJKglue={\hskip 0.10em plus 0.08em minus 0.05em}}%
% —— 纯答案册全线括线（长详解可跨栏断，承练习件全线括线口径；灰底盒不可跨栏断） ——
\ansblockgrayfalse
% —— 双档开关（false 档＝纯值速查：答案印、详解不印；件面开关，不动 sty 本体） ——
\newif\ifansbookdetail
\ifdefined\ansbookdetail
  \ifnum\ansbookdetail=0 \ansbookdetailfalse\else\ansbookdetailtrue\fi
\else
  \ansbookdetailtrue
\fi
% —— 页脚件名（承源件章名；件名＝<件型>答案册） ——
\renewcommand{\qpzhangming}{第二章\quad 平面解析几何}%
\renewcommand{\qpjianming}{导学件答案册}%

\begin{document}
\zhangtitle{第二章\quad 平面解析几何}
\jietitle{2.3.3 直线与圆的位置关系}
\par\glueguard{1}\addvspace{2pt}\noindent\biaoqian{【纯答案册】}{\fangsong 题面不重印\quad 印面号与正册同号同序}\par\addvspace{2pt}

\begin{multicols}{2}
\emergencystretch=1em

\begin{ansblock}[2章-练-课时08-T5]
% ans:2章-练-课时08-T5
\ansitem{1}{A}
\ifansbookdetail
\ansnote{详解}{圆\(x^{2}+(y-1)^{2}=1\)，圆心\((0,1)\)。直线\(y=k(x-1)+1\)过定点\((1,1)\)（\(x=1\)时\(y=1\)），恰为圆心，故\(d=0<r\)，恒相交，故选：A．}
\fi
\end{ansblock}

\begin{ansblock}[2章-练-课时08-T11]
% ans:2章-练-课时08-T11
\ansitem{2}{相切}
\ifansbookdetail
\ansnote{详解}{圆心\((0,0)\)，半径\(r\)。\(d=|-r|/\sqrt{\cos^{2}\theta+\sin^{2}\theta}=r\)，恒相切．}
\fi
\end{ansblock}

\begin{ansblock}[2章-练-课时08-T12]
% ans:2章-练-课时08-T12
\ansitem{3}{[√2, +∞)}
\ifansbookdetail
\ansnote{详解}{表圆：\(\sqrt{m}\)有意义即\(m\geqslant 0\)，\(r^{2}=m^{2}+m-2>0\)，在\(m\geqslant 0\)下\(m>1\)。\par\noindent 与\(x\)轴有公共点\(\iff|\sqrt{m}|\leqslant r\)，即\(m\leqslant m^{2}+m-2\)，即\(m^{2}\geqslant 2\)，得\(m\geqslant\sqrt{2}\)。故\(m\in[\sqrt{2},\ +\infty)\)．}
\fi
\end{ansblock}

\begin{ansblock}[2章-练-课时08-T13]
% ans:2章-练-课时08-T13
\ansitem{4}{(−∞, −5√2/4)∪(5√2/4, +∞)}
\ifansbookdetail
\ansnote{详解}{视线不被挡\(\iff\)直线\(AB\)与圆相离。\(k_{AB}=a/5\)，直线：\(ax-5y+2a=0\)。\par\noindent \(d=|3a|/\sqrt{a^{2}+25}>1\)，即\(9a^{2}>a^{2}+25\)，\(8a^{2}>25\)，\(|a|>5/(2\sqrt{2})=5\sqrt{2}/4\)．}
\fi
\end{ansblock}

\begin{ansblock}[2章-练-课时08-T8]
% ans:2章-练-课时08-T8
\ansitem{5}{C}
\ifansbookdetail
\ansnote{详解}{\(x\geqslant 0\)：\(y=x-1\)，\(x^{2}+(x-1)^{2}=1\)，即\(2x^{2}-2x=0\)，\(x=0\)（\(y=-1\)）或\(x=1\)（\(y=0\)）；\(x<0\)：\(y=-x-1\)，\(x^{2}+(x+1)^{2}=1\)，即\(2x^{2}+2x=0\)，\(x=-1\)（\(y=0\)）（\(x=0\)不在此段）。\par\noindent 交点\((0,-1)\)、\((1,0)\)、\((-1,0)\)共3个，真子集\(2^{3}-1=7\)个，故选：C．}
\fi
\end{ansblock}

\begin{ansblock}[2章-练-课时08-T1]
% ans:2章-练-课时08-T1
\ansitem{6}{A}
\ifansbookdetail
\ansnote{详解}{\((x-1)^{2}+(y-2)^{2}=9\)，圆心\((1,2)\)，\(r=3\)。\(d=|3+8-1|/5=2\)。弦长\(=2\sqrt{9-4}=2\sqrt{5}\)，故选：A．}
\fi
\end{ansblock}

\begin{ansblock}[2章-练-课时08-T2]
% ans:2章-练-课时08-T2
\ansitem{7}{A}
\ifansbookdetail
\ansnote{详解}{过圆内点的弦以垂直于\(OP\)的弦最短。\(k_{OP}=2\)，故所求斜率\(-1/2\)：\(y-2=-(1/2)(x-1)\)，即\(x+2y-5=0\)，故选：A．}
\fi
\end{ansblock}

\begin{ansblock}[2章-练-课时08-T3]
% ans:2章-练-课时08-T3
\ansitem{8}{C}
\ifansbookdetail
\ansnote{详解}{圆心\((1,-2)\)，\(r^{2}=5-m\)。\(d=|1-2-3|/\sqrt{2}=2\sqrt{2}\)。半弦长的平方\(=r^{2}-d^{2}\)，即\(1=5-m-8\)，得\(m=-4\)，故选：C．}
\fi
\end{ansblock}

\begin{ansblock}[2章-练-课时08-T4]
% ans:2章-练-课时08-T4
\ansitem{9}{B}
\ifansbookdetail
\ansnote{详解}{\(|AB|\geqslant 2\sqrt{3}\)，即\(d=\sqrt{r^{2}-(|AB|/2)^{2}}\leqslant\sqrt{4-3}=1\)。又\(d=|m|/\sqrt{2}\)，故\(|m|\leqslant\sqrt{2}\)，故选：B．}
\fi
\end{ansblock}

\begin{ansblock}[2章-练-课时08-T6]
% ans:2章-练-课时08-T6
\ansitem{10}{D}
\ifansbookdetail
\ansnote{详解}{圆\((x-2)^{2}+y^{2}=1\)，圆心\((2,0)\)，\(r=1\)。①过原点（两截距同为0）：\(y=kx\)，\(d=|2k|/\sqrt{1+k^{2}}=1\)，即\(k^{2}=1/3\)，2条；②不过原点：\(x/a+y/a=1\)即\(x+y=a\)，\(d=|2-a|/\sqrt{2}=1\)，\(a=2\pm\sqrt{2}\)（均非0），2条。\par\noindent 共4条，故选：D．}
\fi
\end{ansblock}

\begin{ansblock}[2章-练-课时08-T7]
% ans:2章-练-课时08-T7
\ansitem{11}{B}
\ifansbookdetail
\ansnote{详解}{圆与两轴都切且过\((2,1)\)（第一象限），则圆心\((a,a)\)，\(r=a\)，\(a>0\)。\((2-a)^{2}+(1-a)^{2}=a^{2}\)，即\(a^{2}-6a+5=0\)，\(a=1\)或\(a=5\)。\par\noindent \(a=1\)：圆心\((1,1)\)，\(d=|2-1-3|/\sqrt{5}=2/\sqrt{5}=2\sqrt{5}/5\)；\(a=5\)：圆心\((5,5)\)，\(d=|10-5-3|/\sqrt{5}=2/\sqrt{5}\)，相同。故\(d=2\sqrt{5}/5\)，故选：B．}
\fi
\end{ansblock}

\begin{ansblock}[2章-练-课时08-T9]
% ans:2章-练-课时08-T9
\ansitem{12}{A}
\ifansbookdetail
\ansnote{详解}{圆\((x-2)^{2}+(y+3)^{2}=1\)，圆心\(C(2,-3)\)，\(r=1\)。切线长\(=\sqrt{|PC|^{2}-r^{2}}\)，最小即\(|PC|\)最小\(=C\)到直线距离\(=|4+3-3|/\sqrt{5}=4/\sqrt{5}\)。\par\noindent 最小切线长\(=\sqrt{16/5-1}=\sqrt{11/5}=\sqrt{55}/5\)，故选：A．}
\fi
\end{ansblock}

\begin{ansblock}[2章-练-课时08-T10]
% ans:2章-练-课时08-T10
\ansitem{13}{C}
\ifansbookdetail
\ansnote{详解}{圆\((x-3)^{2}+(y-1)^{2}=9\)，圆心\((3,1)\)，\(r=3\)。对称轴过圆心：\(3+a-1=0\)，得\(a=-2\)，\(A(-1,-2)\)。\(|AC|=\sqrt{16+9}=5\)，切线长\(|AB|=\sqrt{25-9}=4\)，故选：C．}
\fi
\end{ansblock}

\begin{ansblock}[2章-练-课时08-T14]
% ans:2章-练-课时08-T14
\ansitem{14}{√3}
\ifansbookdetail
\ansnote{详解}{设直线\(y=\sqrt{3}x+b\)，\(A(0,b)\)。相切：\(d=|b-1|/2=1\)，得\(b=3\)或\(b=-1\)，\(|AC|=2\)（\(C(0,1)\)）。\(|AB|=\sqrt{|AC|^{2}-r^{2}}=\sqrt{4-1}=\sqrt{3}\)．}
\fi
\end{ansblock}

\begin{ansblock}[2章-练-课时08-T15]
% ans:2章-练-课时08-T15
\ansitem{15}{(1) x=3 或 3x−4y−5=0；(2) a=0 或 a=4/3；(3) a=−3/4}
\ifansbookdetail
\ansnote{详解}{\(C(1,2)\)，\(r=2\)；\(|MC|=\sqrt{4+1}=\sqrt{5}>2\)，\(M\)在圆外。\par\noindent (1) 斜率不存在时\(x=3\)：\(d=|3-1|=2=r\)，是切线。斜率存在时设\(y-1=k(x-3)\)，即\(kx-y-3k+1=0\)：\(d=|k-2-3k+1|/\sqrt{k^{2}+1}=|2k+1|/\sqrt{k^{2}+1}=2\)，即\(4k^{2}+4k+1=4k^{2}+4\)，得\(k=3/4\)，切线\(3x-4y-5=0\)。\par\noindent (2) \(d=|a-2+4|/\sqrt{a^{2}+1}=|a+2|/\sqrt{a^{2}+1}=2\)，即\(a^{2}+4a+4=4a^{2}+4\)，\(3a^{2}-4a=0\)，得\(a=0\)或\(a=4/3\)。\par\noindent (3) 弦长\(2\sqrt{3}\)，半弦\(\sqrt{3}\)，\(d=\sqrt{4-3}=1\)：\(|a+2|/\sqrt{a^{2}+1}=1\)，即\(a^{2}+4a+4=a^{2}+1\)，得\(a=-3/4\)．}
\fi
\end{ansblock}

\begin{ansblock}[2章-练-课时08-T16]
% ans:2章-练-课时08-T16
\ansitem{16}{(1) x²+y²=36；(2) 一定进入，航行时间 1 小时}
\ifansbookdetail
\ansnote{详解}{(1) \(O\)为原点，\(A(3,0)\)，\(B(12,0)\)。由\(|PA|/|PB|=1/2\)：\(4[(x-3)^{2}+y^{2}]=(x-12)^{2}+y^{2}\)，展开：\(4x^{2}-24x+36+4y^{2}=x^{2}-24x+144+y^{2}\)，即\(3x^{2}+3y^{2}=108\)，得\(x^{2}+y^{2}=36\)。\par\noindent (2) \(C(0,-2\sqrt{11})\)。北偏东\(30^{\circ}\)即航线倾角\(60^{\circ}\)，斜率\(\sqrt{3}\)，航线：\(y+2\sqrt{11}=\sqrt{3}x\)，即\(\sqrt{3}x-y-2\sqrt{11}=0\)。\(d=|2\sqrt{11}|/\sqrt{3+1}=\sqrt{11}<6\)，相交，必进入。\par\noindent 弦长\(=2\sqrt{36-11}=10\)，\(t=10/10=1\)（小时）．}
\fi
\end{ansblock}

\begin{ansblock}[2章-导-课时08-G1]
% ans:2章-导-课时08-G1
\ansitem{17}{D}
\ifansbookdetail
\ansnote{详解}{圆心\(O(0,0)\)，\(r=2\)。\(O\)到直线距离\(d=|2|/\sqrt{3+1}=1\)。弦心距\(OM=1\)，\(|AM|=\sqrt{r^{2}-d^{2}}=\sqrt{3}\)，故\(\sin\angle AOM=\sqrt{3}/2\)，\(\angle AOM=60^{\circ}\)，圆心角\(\angle AOB=120^{\circ}\)，故选：D．}
\fi
\end{ansblock}

\begin{ansblock}[2章-导-课时08-G2]
% ans:2章-导-课时08-G2
\ansitem{18}{B}
\ifansbookdetail
\ansnote{详解}{\(d=|-1|/\sqrt{\sin^{2}\theta+1}\)。由\(\theta\neq k\pi+\pi/2\)知\(\sin\theta\neq\pm 1\)，故\(0\leqslant\sin^{2}\theta<1\)，\(d=1/\sqrt{1+\sin^{2}\theta}\in(1/\sqrt{2},\ 1]\)，恒大于\(r=\sqrt{2}/2=1/\sqrt{2}\)，故相离，故选：B．}
\fi
\end{ansblock}

\begin{ansblock}[2章-导-课时08-G3]
% ans:2章-导-课时08-G3
\ansitem{19}{C}
\ifansbookdetail
\ansnote{详解}{\((\sqrt{3})^{2}+(3-2)^{2}=4\)，点在圆上，切线唯一。圆心\((0,2)\)，\(k_{OC}=(3-2)/(\sqrt{3}-0)=\sqrt{3}/3\)，切线斜率\(k=-\sqrt{3}\)，倾斜角\(120^{\circ}\)，故选：C．}
\fi
\end{ansblock}

\begin{ansblock}[2章-导-课时08-G4]
% ans:2章-导-课时08-G4
\ansitem{20}{[−1−2√2, −1+2√2]}
\ifansbookdetail
\ansnote{详解}{圆心\((a,0)\)，\(r=2\)。有公共点\(\iff d=|a+1|/\sqrt{2}\leqslant 2\)，即\(|a+1|\leqslant 2\sqrt{2}\)，得\(-1-2\sqrt{2}\leqslant a\leqslant -1+2\sqrt{2}\)．}
\fi
\end{ansblock}

\begin{ansblock}[2章-导-课时08-G5]
% ans:2章-导-课时08-G5
\ansitem{21}{x−y−3=0}
\ifansbookdetail
\ansnote{详解}{圆心\(C(1,0)\)。\(k_{CP}=(-1-0)/(2-1)=-1\)。\(CP\perp AB\)，故\(k_{AB}=1\)。\(y+1=1\cdot(x-2)\)，即\(x-y-3=0\)．}
\fi
\end{ansblock}

% ---- 尾块（\tailfill[30mm] 定高参——钉档 §三.3：无参在平衡栏呈「内容尾随框」，贴栏底须定高参；实测无参致末页平衡 Overfull\vbox 12.99pt，定高 30mm 三零） ----
\tailfill[30mm]

\end{multicols}
\end{document}
